\section{Introduction}
While multiplication is conceptually simple, implementing it efficiently becomes difficult at large sizes.
Large-integer multiplication is especially important in cryptography, where common sizes include 160--512 bits for Elliptic Curve Cryptography (ECC) and 1024--4096 bits for RSA~\cite{patelComparisonKeySize2023}.
These sizes exceed the native wordlengths of typical hardware resources and must be built from smaller primitives.
% Squaring a large integer is also a common operation in many applications; for example, in RSA, approximately 22.2\% of the operations are squaring operation~\cite{juCompilerLikeFrameworkOptimizing2024}.
% Squaring is the specific instance of multiplication where the multiplier and the multiplicand are identical.
% Although squaring can be performed using general multiplication hardware, specialized squarers exploiting this property can offer significant advantages in terms of resource efficiency.

Field-Programmable Gate Arrays (FPGAs) offer a flexible architecture that can be used to implement a variety of applications in domains such as image processing
~\cite{sarrafzadehImageRestorationFPGA2003}, and cryptography~\cite{prasannaFHEPerformanceModelingFPGA2021}.
Because multiplication is a common application and expensive in LUT fabric, modern FPGAs provide dedicated DSP blocks that deliver higher throughput and better area/energy efficiency than LUT-only implementations.
DSP blocks also support common arithmetic patterns (e.g. MAC), but achieving an optimal balance of DSPs and LUTs for large, structured arithmetic remains challenging.

% FPGAs have two main resources for computing functions, Look-Up Tables (LUTs) and Digital Signal Processing (DSP) blocks.
% While LUT-based architecture provides flexibility, certain functions appear frequently across various applications. 
% Arithmetic operators, for instance, can be inefficient when implemented using LUTs due to significant area and power consumption.
% To address this issue, FPGAs have introduced hard blocks, like DSP blocks, in their architectures, which provide optimized implementations of common functions, making them more efficient than their corresponding LUT-based implementation.
% However, these DSP blocks are a scarse resource compared to LUTs. 
% For example, the Alveo U250 FPGA has only around 12K DSP slices, but has 1728K LUTs as shown in \cite{xilinxU250DataSheet2023}. 
% As a result, applications must make trade-offs between using fast but limited DSP blocks or slow but widespread LUTs.

Large multipliers are typically constructed by decomposing operands into smaller tiles that match the underlying DSP wordlengths.
Schoolbook decomposition forms all partial products and combines them with shifts and additions, while Karatsuba~\cite{karatsubaOriginalPaper} reduces the number of multiplications by trading multipliers for additions (Section~\ref{sec_background}).
For extremely large sizes, FFT~\cite{yapFFTquickmul2000} and NTT~\cite{weemsNTTHighPrecisionInteger2011} based methods are also possible.

However, effective decomposition on FPGAs must be aware of the underlying architecture: DSP primitives are often asymmetric and non-power-of-2, signed/unsigned handling affects packing, and achieving high frequency requires respecting DSP pipeline constraints.
FPGAs also expose structure such as cascade paths that can improve both performance and resource usage when exploited.
This paper addresses large integer multiplication in an architecture-aware manner, enabling both efficient base multipliers and scalable large multipliers on FPGAs.

The main contributions of this paper are: 
\begin{enumerate}
    \item An architecture-aware base multiplier that maps efficiently to FPGA DSP blocks while achieving high frequency and low resource usage.
    \item A comparison against a vendor $64\times64$ IP multiplier.
    \item Automated construction of large multipliers from the proposed base design using Schoolbook and Karatsuba decompositions, with FPGA-aware pipelining and mapping.
    \item Evaluation across operand sizes and base sizes, including a $1024\times1024$-bit design compared against an HLS implementation and the approach in~\cite{zhangIMpressLargeInteger2022} and~.
\end{enumerate}

Overall, our approach achieves up to 56\% higher frequency while using about 17\% fewer DSP blocks compared to prior approaches.

